\documentclass[conference]{IEEEtran}
\IEEEoverridecommandlockouts

\usepackage{cite}
\usepackage{amsmath,amssymb,amsfonts}
\usepackage{algorithmic}
\usepackage{graphicx}
\usepackage{textcomp}
\usepackage{xcolor}
\usepackage{booktabs}
\usepackage{url}
\usepackage{multirow}
\usepackage{microtype}

\begin{document}

\title{PathogenIQ: Machine Learning for Metagenomic Wastewater Surveillance}

\author{
  \IEEEauthorblockN{Richmond Azumah\textsuperscript{1}}
  \IEEEauthorblockA{\textsuperscript{1}Department of Computer Science,
    Georgia State University, Atlanta, GA\\
    razumah1@student.gsu.edu}
}

\maketitle

% ── Abstract ──────────────────────────────────────────────────────────────────
\begin{abstract}
Metagenomic sequencing-based biosurveillance provides population-level
pathogen monitoring without individual testing, but existing bioinformatics
pipelines lack integrated machine learning for risk stratification, temporal
outbreak detection, and automated alerting. We present \textit{PathogenIQ},
a modular, open-source, end-to-end pipeline that ingests Kraken2 taxonomic
reports and integrates: (1) a false-discovery-rate (FDR)-corrected Spearman
co-occurrence graph with stochastic block model (SBM) community detection;
(2) a multi-signal composite risk scorer combining curated pathogen
abundance weights, community context, and isolation-forest-based
cohort-relative anomaly (novelty) detection; (3) cumulative sum (CUSUM)
changepoint and Mann-Kendall trend analysis; and (4) a real-time web
dashboard with Slack/email alerting.
On 60 real wastewater samples ($\sim$30-week site-207 series),
PathogenIQ classifies 55 runs as MODERATE and 5 as HIGH, where an
abundance-only threshold raises none; we show this divergence reflects
PathogenIQ's operational decision logic (a fixed community-context
offset and a cohort-recalibrated direct-detection rule supplying the
threshold margin) rather than superior raw detection. The CUSUM layer
transiently crosses its decision boundary ($C_t=4.54$ vs.\ $h=4.0$) around
week 20 before receding to baseline, capturing a genuine risk excursion.
On six controlled synthetic scenarios ($N=52$), the abundance-driven
pathway reaches Sensitivity=98.8\%, F1=99.3\% with zero false positives
(FAR=0.0\% across 10 seeds), a 50-point FAR reduction over a na\"{i}ve
threshold, while community and novelty signals add interpretable context
without changing binary alert decisions. PathogenIQ's central contribution
is this end-to-end integration, from raw Kraken2 reports to a deployable
alerting dashboard.
\end{abstract}

\begin{IEEEkeywords}
biosurveillance, metagenomics, stochastic block model, CUSUM,
wastewater epidemiology, pathogen risk scoring, applied machine learning,
public health signal processing
\end{IEEEkeywords}

% ── 1. Introduction ───────────────────────────────────────────────────────────
\section{Introduction}
\label{sec:intro}

Wastewater-based epidemiology (WBE) provides population-level pathogen
surveillance without individual clinical testing, detecting pathogen
signals days to weeks before hospital admissions or case counts
rise~\cite{Bibby2021,Wu2020,Peccia2020,Medema2020}. During the COVID-19
pandemic, SARS-CoV-2 RNA concentrations in sewage tracked community
prevalence with a 4--7 day lead over reported cases~\cite{Daughton2020},
and genome-resolved wastewater metagenomics enabled detection of
regionally prevalent variants~\cite{CritsChristoph2021}. The same
metagenomic approach extends well beyond respiratory viruses: large-scale
wastewater sequencing now supports global antimicrobial resistance (AMR)
monitoring~\cite{Hendriksen2019} and broader multi-pathogen
surveillance~\cite{Quince2017}.

These applications typically rely on Kraken2 with Bracken
re-estimation~\cite{Wood2019,Lu2017} (see~\cite{Breitwieser2019} for a
review) to produce per-sample taxonomic abundance reports. These reports
are accurate and information-rich, but the analytical layer that would turn
them into public-health decisions is missing: practitioners must manually
inspect hundreds of taxa per report, apply informal abundance thresholds,
and track trends across time by hand. At scale (a metropolitan program
monitoring dozens of sites accumulates tens of thousands of reports
annually, each requiring expert review), this bottleneck introduces
days of latency and makes continuous surveillance economically infeasible
without automation.

Machine learning offers principled tools to close this gap, and the
signal-processing perspective of MLSP (changepoint detection, graph
signal analysis, and anomaly scoring) maps directly onto the WBE problem.
Yet prior work has addressed these tools in isolation. Deep
learning has been applied to metagenomic classification~\cite{Liang2020},
and reference-free variant detection in sequencing reads is being explored
in concurrent work~\cite{Azumah2026}; statistical outbreak detection has a
long history in public health signal processing~\cite{Farrington1996,Unkel2012}.
Anomaly detection methods~\cite{Chandola2009} have been applied to
surveillance streams, but not combined with community-level pathogen
co-occurrence signals. No unified, open-source system integrates these
capabilities into a single deployable pipeline on standard Kraken2 outputs.

We present \textbf{PathogenIQ}, which contributes:
\begin{itemize}
  \item an open-source, end-to-end ML pipeline from raw Kraken2 reports to
    interpretable risk scores, unifying community detection, multi-signal
    risk scoring, temporal anomaly detection, and automated alerting in one
    deployable system;
  \item an FDR-controlled Spearman co-occurrence graph partitioned by a
    stochastic block model (SBM) to expose pathogen co-colonization
    structure invisible to individual taxon abundances;
  \item a three-signal composite risk scorer (150-taxon abundance database,
    SBM community context, isolation-forest novel-pathogen detection) on a
    calibrated $[0,1]$ scale; and
  \item CUSUM changepoint and Mann-Kendall trend analysis over a per-site
    history, paired with a dashboard and Slack/email alerting for early
    warning without analyst intervention.
\end{itemize}

% ── 2. Related Work ───────────────────────────────────────────────────────────
\section{Related Work}

\textbf{Wastewater surveillance pipelines.}
Standard WBE pipelines (Freyja~\cite{Karthikeyan2022}, COJAC~\cite{Jahn2022})
focus on SARS-CoV-2 variant deconvolution and are reference-dependent and
single-pathogen. Broader genomic surveillance~\cite{CritsChristoph2021,Ahmed2020}
and global AMR monitoring~\cite{Hendriksen2019} demonstrate multi-pathogen
potential but lack quantitative ML-based risk scoring.
Recent WBE systems~\cite{Msomi2025,Gauthier2025,Kohle2024,Feistel2025,Pagsuyoin2025}
have begun integrating dashboards, ML forecasting, and metagenomic
bacterial detection, yet none unifies all capabilities in one pipeline
(Table~\ref{tab:comparison}).

\textbf{Outbreak detection in public health.}
Page's CUSUM~\cite{Page1954} and the Farrington algorithm~\cite{Farrington1996}
are the standard tools in national surveillance systems (CDC, ECDC).
Unkel et al.~\cite{Unkel2012} review prospective surveillance algorithms;
most assume discrete case-count data rather than continuous ML-derived
risk scores. PathogenIQ adapts CUSUM to per-site metagenomics risk scores,
enabling continuous monitoring without Poisson count assumptions.

\textbf{Community detection in microbial ecology.}
The SBM~\cite{Holland1983,Karrer2011} has been widely applied to complex
networks~\cite{Newman2006,Blondel2008}. In microbiology, Spearman-based
co-occurrence networks~\cite{Barber2012} and compositionally robust
inference~\cite{Kurtz2015} characterize co-colonization; Faust \&
Raes~\cite{Faust2012} survey microbial interaction models broadly.
PathogenIQ is, to our knowledge, the first pathogen surveillance pipeline
to incorporate SBM community detection over an FDR-corrected Spearman
co-occurrence network as a structural component of its composite
risk-scoring framework.

\textbf{Automated surveillance dashboards.}
Tools such as Nextstrain~\cite{Hadfield2018} demonstrate the value of
real-time dashboards for genomic epidemiology. PathogenIQ extends this
paradigm to the metagenomics setting, coupling ML risk scores directly
to a deployable alerting interface.

\begin{table}[h]
\centering
\caption{Capability comparison of recent WBE computational surveillance
  systems, framed by practitioner-facing benefit.
  \textcolor{green!60!black}{\checkmark} = supported;
  \textcolor{red}{$\times$} = not supported.}
\label{tab:comparison}
\resizebox{\columnwidth}{!}{%
\begin{tabular}{lccccccc}
\toprule
\textbf{System} &
  \rotatebox{60}{\textbf{Tracks many pathogens}} &
  \rotatebox{60}{\textbf{Automated risk score}} &
  \rotatebox{60}{\textbf{Co-colonization risk}} &
  \rotatebox{60}{\textbf{Early-warning trends}} &
  \rotatebox{60}{\textbf{Flags novel pathogens}} &
  \rotatebox{60}{\textbf{Automated alerts}} &
  \rotatebox{60}{\textbf{Live dashboard}} \\
\midrule
Msomi et al.~\cite{Msomi2025}         & \textcolor{red}{$\times$}            & \textcolor{red}{$\times$}            & \textcolor{red}{$\times$} & \textcolor{red}{$\times$}            & \textcolor{red}{$\times$} & \textcolor{red}{$\times$} & \textcolor{green!60!black}{\checkmark} \\
Gauthier et al.~\cite{Gauthier2025}   & \textcolor{red}{$\times$}            & \textcolor{red}{$\times$}            & \textcolor{red}{$\times$} & \textcolor{green!60!black}{\checkmark} & \textcolor{red}{$\times$} & \textcolor{red}{$\times$} & \textcolor{red}{$\times$}            \\
Kohle et al.~\cite{Kohle2024}         & \textcolor{red}{$\times$}            & \textcolor{red}{$\times$}            & \textcolor{red}{$\times$} & \textcolor{red}{$\times$}            & \textcolor{red}{$\times$} & \textcolor{red}{$\times$} & \textcolor{red}{$\times$}            \\
Feistel et al.~\cite{Feistel2025}     & \textcolor{red}{$\times$}            & \textcolor{red}{$\times$}            & \textcolor{red}{$\times$} & \textcolor{red}{$\times$}            & \textcolor{red}{$\times$} & \textcolor{red}{$\times$} & \textcolor{red}{$\times$}            \\
Pagsuyoin et al.~\cite{Pagsuyoin2025} & \textcolor{red}{$\times$}            & \textcolor{green!60!black}{\checkmark} & \textcolor{red}{$\times$} & \textcolor{green!60!black}{\checkmark} & \textcolor{red}{$\times$} & \textcolor{red}{$\times$} & \textcolor{red}{$\times$}            \\
\textbf{PathogenIQ (ours)}            & \textcolor{green!60!black}{\checkmark} & \textcolor{green!60!black}{\checkmark} & \textcolor{green!60!black}{\checkmark} & \textcolor{green!60!black}{\checkmark} & \textcolor{green!60!black}{\checkmark} & \textcolor{green!60!black}{\checkmark} & \textcolor{green!60!black}{\checkmark} \\
\bottomrule
\end{tabular}%
}
\end{table}

% ── 3. Methods ────────────────────────────────────────────────────────────────
\section{Methods}

PathogenIQ takes a cohort of Kraken2 \texttt{.report} files as input and
returns per-sample risk scores with automated alerts. As shown in
Fig.~\ref{fig:pipeline}, it comprises six stages: (1) co-occurrence graph
construction, (2) SBM community detection, and (3) isolation-forest
novel-pathogen detection, which together with curated abundance weights
feed (4) multi-signal risk scoring, followed by (5) CUSUM changepoint and
trend analysis and (6) dashboard alerting. We describe each below.

\begin{figure}[t]
  \centering
  \includegraphics[width=\columnwidth]{figures/fig1_pipeline}
  \caption{PathogenIQ pipeline. \textit{Top (preprocessing):}
  wastewater samples are sequenced, classified by Kraken2, and written
  as \texttt{.report} files; only the report feeds PathogenIQ.
  \textit{Bottom:} six numbered stages: (1) co-occurrence graph,
  (2) SBM community detection, and (3) novel-pathogen detection feed
  (4) multi-signal risk scoring, then (5) CUSUM temporal analysis and
  (6) the alerting dashboard.}
  \label{fig:pipeline}
\end{figure}

\subsection{Data Ingestion and Filtering}

PathogenIQ accepts a directory of Kraken2 \texttt{.report} files or a
pre-computed taxa $\times$ samples count matrix. Each report is parsed and
read counts are aggregated at a user-specified taxonomic rank ($G$=genus,
$S$=species, or $F$=family) into a full taxa $\times$ samples count matrix.
This matrix is filtered to retain only taxa present in at least
$\rho_{\min}=10\%$ of samples and with at least $r_{\min}=50$ total reads,
yielding the filtered count submatrix
$\mathbf{X} \in \mathbb{R}^{T \times N}$ of $T$ taxa and $N$ samples that
feeds the co-occurrence graph described next. The $10\%$ prevalence
threshold requires a retained taxon to appear in at least 6 of $N=60$
samples, comfortably above the minimum needed for a non-degenerate
Spearman rank correlation, while $r_{\min}=50$ excludes taxa with too few reads to estimate relative
abundance reliably. Relative abundances are computed column-wise from
$\mathbf{X}$ and missing values are imputed with zero.

\subsection{Co-occurrence Graph Construction}

We build a co-occurrence graph $G = (V, E)$ with edge weights as Spearman
rank correlations $\rho_{ij}$ between taxa across samples.
Benjamini-Hochberg correction~\cite{Benjamini1995} at $\alpha_{\text{FDR}}=0.05$
controls the FDR from $\binom{T}{2}$ tests; an edge is retained only if
significant and $|\rho_{ij}| \geq 0.45$. The resulting adjacency matrix
$\mathbf{A}$ is passed to the SBM.

\subsection{Stochastic Block Model Community Detection}

We train an SBM~\cite{Holland1983,Karrer2011} on $\mathbf{A}$ using
variational EM with ICL model selection~\cite{Peixoto2014}.
The model posits that taxon $i$ belongs to community
$z_i \in \{1,\ldots,K\}$ with mixing proportion $\pi_k = P(z_i=k)$,
and that edges are drawn independently given community assignments:
%
\begin{equation}
P(A_{ij} = 1 \mid z_i = k, z_j = l) = B_{kl}
\end{equation}
%
where $A_{ij}$ is the $(i,j)$ entry of the adjacency matrix $\mathbf{A}$,
$k$ and $l$ are community indices, and
$\mathbf{B} \in [0,1]^{K \times K}$ is the block connection matrix
whose entry $B_{kl}$ gives the edge probability between communities
$k$ and $l$. Soft membership assignments
$q_{ik} = P(z_i = k \mid \mathbf{A})$ are updated via:
%
\begin{equation}
\log q_{ik} \!\propto\! \log \pi_k + \!\sum_{j \neq i}\! \sum_l q_{jl}
  \big[ A_{ij} \log B_{kl} + (1 \!-\! A_{ij}) \log(1 \!-\! B_{kl}) \big]
\end{equation}
%
Model order $K$ is selected by maximizing ICL, with $n_{\text{init}}=10$
restarts and $K_{\max}=\min(12, T-1)$.
A community $k$ is flagged \emph{pathogen-enriched} if $\geq 20\%$ of
its members are in the curated pathogen database.
Partition stability is assessed via $B=100$ bootstrap resamples
(Section~\ref{sec:results}).

\subsection{Multi-Signal Risk Scoring}
\label{sec:recalibration}

For each sample $n$, PathogenIQ combines three signals into a composite
risk score $\hat{r}_n \in [0,1]$: a pathogen abundance score
$s_n^{(\text{ab})}$, a community context flag $s_n^{(\text{comm})}$,
and a novel-pathogen score $s_n^{(\text{nov})}$:
%
\begin{equation}
\hat{r}_n = \max\!\left(
  \alpha s_n^{(\text{ab})} + \beta s_n^{(\text{comm})} + \gamma s_n^{(\text{nov})},\;
  d_n
\right)
\label{eq:risk}
\end{equation}
%
where $\alpha=0.50$, $\beta=0.25$, $\gamma=0.25$ are signal weights,
and $d_n$ is a direct-detection score that bypasses the weighted sum
for epidemiologically significant detections (e.g., any BSL-3 genus
at $\geq 5\%$ relative abundance immediately yields $d_n \geq 0.85$).
\textbf{Signal 1 (Pathogen abundance)} $s_n^{(\text{ab})}$:
Each taxon genus $g$ is assigned a curated risk weight $w_g \in [0,1]$
from our database of 150 taxa (bacteria, viruses, fungi, parasites).
Weights were assigned by mapping each taxon to WHO priority pathogen
tiers~\cite{WHO2024} and CDC Select Agent classifications, with BSL-3/4
organisms receiving $w_g \geq 0.8$ and environmental commensals
receiving $w_g \leq 0.2$.
The abundance score is $s_n^{(\text{ab})} = \min(1, \sum_{g} w_g x_{gn})$.

\textbf{Signal 2 (Community context)} $s_n^{(\text{comm})}$:
Binary flag equal to 1 if at least $\phi=15\%$ of sample $n$'s present
taxa belong to a pathogen-enriched SBM community as defined above, and 0
otherwise. This signal captures co-colonisation risk that is invisible to
single-taxon abundance scoring.

\textbf{Signal 3 (Novel Pathogen Detection)} $s_n^{(\text{nov})}$:
An Isolation Forest~\cite{Liu2008,Chandola2009} is trained on $\mathbf{X}$
in a \emph{transductive} setting so that all cohort samples are seen
simultaneously, which is standard practice for cohort-level novel
pathogen detection when no external reference distribution is available.
The normalized anomaly score is clipped to $[0,1]$ and reflects
deviation from the cohort mean profile; it cannot be interpreted as
out-of-sample generalization without an independent reference cohort.

The weights $\alpha=0.50$, $\beta=0.25$, $\gamma=0.25$ are
domain-informed priors: abundance is primary because direct pathogen
presence is the most reliable indicator, while community context and
novel-pathogen detection serve as secondary modifiers. A sensitivity analysis across
four weight configurations (equal: 0.33/0.33/0.33; abundance-heavy:
0.70/0.15/0.15; current: 0.50/0.25/0.25; community-heavy:
0.40/0.40/0.20) produces identical F1=1.000 on all high-signal
scenarios, confirming that performance is insensitive to exact weight
assignment. Risk levels are: LOW $[0, 0.3)$, MODERATE $[0.3, 0.6)$,
HIGH $[0.6, 0.8)$, CRITICAL $[0.8, 1.0]$.

\textbf{Baseline-relative recalibration of $d_n$.} The absolute
thresholds defining $d_n$, including the multi-pathogen
additive-load rule (\texttt{load}~$\geq 0.10 \Rightarrow d_n \geq
0.62$, where \texttt{load} is the sum of risk-weight-weighted relative
abundances of the top detected pathogens), are calibrated against
outbreak-level contamination. On endemic, real-world wastewater
communities, common enteric/environmental genera (e.g.\
\textit{Pseudomonas}, \textit{Streptococcus}) routinely push
\texttt{load} above 0.10 even when nothing unusual is occurring at the
site, causing $d_n$ to saturate and override the weighted composite in
Eq.~\ref{eq:risk} on nearly every sample. We therefore introduce an
optional, post-hoc recalibration: for a cohort of samples from the same
site, we compute the mean $\mu_L$ and standard deviation $\sigma_L$ of
\texttt{load} across the cohort and retain the direct-detection
override only for samples whose load is a statistical anomaly relative
to that baseline, $z_L=(\texttt{load}-\mu_L)/\sigma_L \geq 2$;
otherwise $d_n$ is suppressed and $\hat{r}_n$ falls back to the
weighted composite alone. This recalibration is applied only to the
real-data case study (Section~\ref{sec:real-data}); the synthetic
benchmark (Table~\ref{tab:main}) continues to use
the fixed absolute thresholds, since the injected outbreak scenarios
are designed to exceed any plausible cohort baseline.

\subsection{Temporal Changepoint and Trend Analysis}

PathogenIQ maintains a per-site SQLite history database. At each run,
the risk score $\hat{r}_t$ is appended and the one-sided CUSUM
statistic $C_t$~\cite{Page1954} is updated as:
%
\begin{equation}
C_t = \max(0,\; C_{t-1} + z_t - k)
\label{eq:cusum}
\end{equation}
%
where $C_{t-1}$ is the statistic from the previous run,
$z_t = (\hat{r}_t - \mu_w)/\sigma_w$ is the z-score of the current
risk score relative to the site-specific rolling mean $\mu_w$ and
standard deviation $\sigma_w$ over a window of $W$ weeks,
$k=0.5$ is the allowance parameter that controls sensitivity, and
an alert fires when $C_t \geq h=4.0$.
Because the CUSUM operates on z-score-normalized inputs, $h=4.0$ is
transferable from discrete case-count surveillance~\cite{Farrington1996}
to continuous ML-derived scores without recalibration.
Mann-Kendall $\tau$~\cite{Mann1945,Kendall1975} detects monotonic trends;
Sen's slope~\cite{Sen1968} provides a one-step-ahead forecast.

\subsection{Dashboard and Alerting}

PathogenIQ serves a FastAPI/Uvicorn web dashboard presenting per-site
risk timelines, the co-occurrence network, and per-sample signal
breakdowns. When a CUSUM alert fires or a sample exceeds $\theta=0.6$,
the system dispatches push notifications via Slack webhook or SMTP email,
delivering actionable alerts directly to practitioners without analyst
intervention.

% ── 4. Experimental Setup ─────────────────────────────────────────────────────
\section{Experimental Setup}

\subsection{Dataset}

We evaluate PathogenIQ on a single-facility, time-ordered series of 60
Kraken2 genus-level metagenomic reports drawn from NCBI BioProject
PRJNA957477 (SRA study SRP433461; Verily Life Sciences national
wastewater surveillance program, 14{,}418 runs across many facilities).
We restrict to \emph{Site 207}, a wastewater treatment facility serving
1.25M residents in Nevada (24h composite samples, post-grit-removal). Of
the 281 Site~207 runs in the BioProject, we keep only the 173 sequenced
on an Illumina NovaSeq 6000 (excluding 108 earlier NextSeq~2000 amplicon
runs from 2023-09--2024-09, which are too shallow for genus-level
classification), sort by collection date, and take the most recent
$N=60$ runs (accessions SRR36005439--SRR39112987), spanning
2025-11-05--2026-06-03 ($\sim$30 weeks, sampled approximately twice
weekly). This deterministic selection rule and the full accession list
are provided with our open-source code release.

Each run was classified with Kraken2~v2.1.6 (standard RefSeq database,
$k=35$, built 2025-10-15) using \texttt{--confidence 0.05 --paired} on a
12-thread node; per-run wall-clock time averaged 280s (range 91--498s),
so the 60-run series completes in $\sim$4.7\,h. Clinical ground-truth
labels are unavailable for this cohort, so we present the real-data run
as a case study of system behaviour, with quantitative performance
reported on the controlled synthetic benchmark below.

For controlled validation we evaluate on six synthetic scenarios:
\emph{negative control} ($n=12$), \emph{low contamination} ($n=10$),
\emph{high contamination} ($n=8$), \emph{critical pathogen spike}
($n=6$), \emph{multi-pathogen community} ($n=8$), and \emph{novel taxon
intrusion} ($n=8$); $N=52$ samples total. PathogenIQ contains no
learned parameters trained on these scenarios: signal weights are
fixed domain priors and the pathogen database is sourced from WHO/CDC;
the benchmark is therefore not evaluation on training data.

\subsection{Baselines}

\begin{itemize}
  \item \textbf{Threshold}: Flag any sample where any pathogen-DB genus
    exceeds 1\% relative abundance.
  \item \textbf{Abundance-only}: Signal $s_n^{(\text{ab})}$ alone
    ($\beta=\gamma=0$, $\alpha=1$).
  \item \textbf{SBM-only}: Community flag only, without direct detection
    or novel-pathogen detection.
\end{itemize}

\subsection{Metrics}

We report sensitivity, specificity, F1, and false alert rate (FAR) at
$\theta=0.6$. For temporal detection we report whether CUSUM detects
signals that the single-point threshold baseline misses entirely.

% ── 5. Results ────────────────────────────────────────────────────────────────
\section{Results}
\label{sec:results}

\subsection{Detection Performance}

Table~\ref{tab:main} summarizes synthetic benchmark performance (Poisson
sequencing-depth noise, mean over 10 seeds). PathogenIQ achieves
near-perfect detection (Sensitivity=98.8\% [95.8, 100.0], F1=99.3\%
[97.8, 100.0]) on high-signal scenarios with zero false positives
(FAR=0.0\%, Spec.=100\%), versus FAR=50\% for the Threshold baseline. The
SBM-only baseline raises zero alerts at $\theta=0.6$ (F1 undefined)
because per-scenario sample counts (6--12) are too few for the community
step to stabilize. PathogenIQ's full pipeline matches Abundance-only
\emph{exactly} on every synthetic metric: the community and novelty
signals add context without changing the binary decision here; their
value emerges on the real series (Section~\ref{sec:real-data}), where the
fixed community contribution and a recalibrated direct-detection override
together produce 5 alerts an Abundance-only baseline misses entirely.

\begin{table}[h]
\centering
\caption{Detection performance on synthetic benchmark ($N=52$ samples,
  $\theta=0.6$), with sequencing-depth (Poisson) noise. F1 and Sens.\
  are mean [95\% CI] over 10 seeds.}
\label{tab:main}
\resizebox{\columnwidth}{!}{%
\begin{tabular}{lcccc}
\toprule
\textbf{Method} & \textbf{Sens.} & \textbf{Spec.} & \textbf{F1} & \textbf{FAR} \\
\midrule
Threshold      & 100.0\% & 50.0\% & 100.0\% & 50.0\% \\
Abundance-only & 98.8\% [95.8, 100.0] & 100.0\% & 99.3\% [97.8, 100.0] & 0.0\% \\
SBM-only       & 0.0\%  & 100.0\% & n/a$^*$   & 0.0\% \\
PathogenIQ     & 98.8\% [95.8, 100.0] & 100.0\% & 99.3\% [97.8, 100.0] & 0.0\% \\
\bottomrule
\multicolumn{5}{l}{\footnotesize $^*$Undefined: SBM-only raises 0 alerts (precision $0/0$).}\\
\end{tabular}%
}
\end{table}

\subsection{Synthetic Benchmark Validation}

Per-scenario, PathogenIQ matches Abundance-only \emph{exactly} (Poisson
noise, 10 seeds): F1=1.000 on the high-contamination and critical-spike
scenarios, F1=0.980 [0.933, 1.000] on the noisy multi-pathogen community,
specificity 1.00 on the negative and low-contamination controls, and F1
undefined by design on novel-taxon intrusion (sub-threshold for all
methods). At $\theta=0.6$ the abundance pathway dominates whenever a
pathogen is present at benchmark concentrations, so the community and
novelty terms act as \emph{secondary modifiers} that add context without
moving F1; their value lies in the real-data case study
(Section~\ref{sec:real-data}), not in synthetic F1.

To delimit exactly when the community signal \emph{can} change a decision
on independently drawn samples, we sweep the per-pathogen abundance of a
co-occurring four-pathogen group (\emph{Salmonella}, \emph{Escherichia},
\emph{Listeria}, \emph{Shigella}; correlation 0.9; $1$--$8\%$, 10 seeds)
and find its marginal effect on the alert rate is $0.00$ at every level:
below $\sim$$4\%$ neither method alerts, and above it the direct-detection
load rule fires for both. This is structural: the community and novelty
terms contribute at most $\beta+\gamma=0.5$, so they can lift a sample
across $\theta=0.6$ only once abundance already reaches the
direct-detection band the baseline also triggers. On independently drawn
samples the two context signals are thus redundant with abundance by
construction; their decision-relevance emerges only under the cohort
recalibration of Section~\ref{sec:recalibration}.

\begin{figure}[h]
  \centering
  \includegraphics[width=\columnwidth]{figures/fig3_sbm_network}
  \caption{$K=2$ SBM partition of the real-data co-occurrence network
    ($N$=60 samples, 54 taxa). Force-directed layout of the 29/54 taxa
    with $\geq$1 FDR-significant Spearman edge ($|\rho|\geq0.45$,
    BH-corrected $\alpha$=0.05; 28 of 1431 possible pairs); the
    remaining 25 taxa have no significant edges and are omitted.
    Community 1 ($n=4$: \emph{Escherichia}, \emph{Enterobacter},
    \emph{Citrobacter}, \emph{Vibrio}, the enteric/fecal-indicator
    cluster, highlighted) is fully embedded in the largest connected
    component (19 taxa, including 15 Community-0 members); the
    remaining 10 connected Community-0 taxa form four small
    components, e.g.\ Pseudomonas (38.9\% mean abundance, the most
    abundant genus overall), Microbacterium (26.4\%), Chryseobacterium
    (16.7\%), and Ralstonia. Node size $\propto \sqrt{\text{mean
    abundance}}$.}
  \label{fig:sbm_network}
\end{figure}

\subsection{Temporal Detection: CUSUM Dynamics}

Fig.~\ref{fig:cusum} shows the CUSUM trajectory over the 60-run,
$\sim$30-week site-207 series (recalibrated scores,
Section~\ref{sec:real-data}). Risk scores fluctuate between 0.375 and
0.622 (mean 0.499, $\sigma$=0.063), and five runs cross $\theta=0.6$
into the HIGH band. The first two occur in close succession (weeks 18
and 20), and $C_t$ accumulates across this cluster to a series maximum
of $4.54$ at week 20 (the only point satisfying the alert criterion
$C_t \geq h=4.0$, where a real-time deployment would have alerted)
before receding to $0.81$ by series end as runs return toward baseline;
the other three HIGH runs occur in isolation and produce no comparable
accumulation. This complements the synthetic benchmark
(Table~\ref{tab:main}): the high-signal scenarios show the
direct-detection pathway fires reliably when warranted (Sens.=98.8\%),
while the real series shows the temporal layer responds to a genuine
cluster and then returns to quiescence rather than over-alerting. Since
the synthetic scenarios are independently drawn with no temporal
ordering, the real series is the only evidence for CUSUM's behaviour
under realistic dynamics, now including both an excursion above $h$ and
recovery to baseline.

\begin{figure}[h]
  \centering
  \includegraphics[width=\columnwidth]{figures/fig2_cusum}
  \caption{CUSUM trace over the real $\sim$30-week, 60-run site-207
    series (recalibrated risk scores). Top: per-run risk score (dashed
    line: alert threshold $\theta=0.6$; five runs transiently exceed
    it). Bottom: CUSUM statistic $C_t$ (dashed line: decision boundary
    $h=4.0$). $C_t$ rises to a series maximum of 4.54 at week 20
    (exceeding $h$) before receding to 0.81 by series end,
    demonstrating both sensitivity to a genuine cluster of elevated
    runs and recovery to baseline thereafter.}
  \label{fig:cusum}
\end{figure}

\subsection{Real-Data Signal Breakdown}
\label{sec:real-data}

Fig.~\ref{fig:risk_dist} shows the three-signal decomposition across
all 60 real wastewater samples after the recalibration of
Section~\ref{sec:recalibration}. The recalibrated direct-detection
override $d_n$ remains active for 2 of 60 samples (2026-01-14 and
2026-04-01; cohort load $z$-scores of 2.25 and 2.03, both
$\geq z_L=2$), so for these two $\hat{r}_n=d_n=0.62$; for the
remaining 58 samples $\hat{r}_n$ equals the weighted composite
$0.5\,s_n^{(\text{ab})}+0.25\,s_n^{(\text{comm})}+0.25\,s_n^{(\text{nov})}$,
which ranges 0.375--0.622 (mean 0.497, $\sigma$=0.060) across the
cohort (55 MODERATE, 5 HIGH overall).

The abundance signal $s_n^{(\text{ab})}$ varies continuously
(0.112--0.500, mean 0.283) and the novel-pathogen signal
$s_n^{(\text{nov})}$ varies even more widely (0.078--0.923, mean
0.423); together these two continuous signals account for essentially
all of the sample-to-sample variation visible in
Fig.~\ref{fig:risk_dist}. The community signal $s_n^{(\text{comm})}$
($\beta=0.25$), by contrast, is active (=1) in \emph{all} 60 of 60
samples: at least $\phi=15\%$ of every sample's present taxa fall in
the pathogen-enriched SBM community throughout the $\sim$30-week
series. As on the synthetic benchmark (Section~\ref{sec:results}), a
constant signal cannot discriminate between samples; on this real
cohort, however, its fixed $+0.25$ contribution is the deciding factor
for 3 of the 5 HIGH classifications (2026-01-07, 2026-03-16,
2026-05-04): for these runs $0.5\,s_n^{(\text{ab})}+0.25\,s_n^{(\text{nov})}$
alone is only 0.362--0.372 (MODERATE), but adding the constant community
term pushes $\hat{r}_n$ to 0.612--0.622, crossing $\theta=0.6$ into
HIGH. The other 2 HIGH classifications (2026-01-14, 2026-04-01) arise
from the recalibrated direct-detection override, independent of the
community term. Table~\ref{tab:real-ablation} quantifies the divergence
from the Abundance-only baseline: because $s_n^{(\text{ab})}$ never
exceeds 0.500 on this cohort, Abundance-only raises zero alerts across
all 60 runs, while PathogenIQ raises 5.

\begin{table}[h]
\centering
\caption{Real-data comparison on the 60-run site-207 series:
  risk-level distribution and alert count ($\hat{r}_n \geq
  \theta=0.6$) for PathogenIQ vs.\ the Abundance-only baseline
  ($s_n^{(\text{ab})}$ alone, $\alpha=1,\beta=\gamma=0$).}
\label{tab:real-ablation}
\begin{tabular}{lcccc}
\toprule
\textbf{Method} & \textbf{LOW} & \textbf{MODERATE} & \textbf{HIGH} & \textbf{Alerts} \\
\midrule
Abundance-only    & 34 & 26 & 0 & 0 \\
PathogenIQ (Full) & 0  & 55 & 5 & 5 \\
\bottomrule
\end{tabular}
\end{table}

\begin{figure}[h]
  \centering
  \includegraphics[width=\columnwidth]{figures/fig4_risk_dist}
  \caption{Risk signal distributions across 60 real wastewater samples
    (recalibrated). Abundance, novel-pathogen, and composite scores
    shown as violins (continuous); composite equals the weighted sum
    except for 2 samples with an active recalibrated direct-detection
    override (Section~\ref{sec:real-data}). Community signal shown as
    stacked bar (binary): 60/60 samples (100\%) active, contributing a
    constant $\beta=0.25$ offset that is nonetheless the deciding
    factor for 3 of the 5 HIGH classifications.}
  \label{fig:risk_dist}
\end{figure}

\subsection{Ablation: Signal Contribution}

A leave-one-out ablation on the synthetic benchmark (single seed,
$\theta=0.6$) confirms this regime structure: removing the community or
novel-pathogen signal leaves F1 unchanged ($\Delta$F1$=0.000$ each),
whereas removing
the direct-detection pathway collapses it ($\Delta$F1$=-1.000$),
confirming it is the dominant signal for clear outbreak events. On the
real series the balance shifts: with the override active for only 2 of
60 samples, the continuous abundance and novelty signals drive the
magnitude of $\hat{r}_n$ and the CUSUM dynamics, while the constant
community flag supplies the threshold margin for the remaining alerts
(Section~\ref{sec:real-data}).

% ── 6. Discussion and Conclusion ──────────────────────────────────────────────
\section{Discussion and Conclusion}

PathogenIQ's central contribution is the integration of FDR-corrected
co-occurrence graph construction, SBM community detection, multi-signal
composite risk scoring, temporal changepoint monitoring, and an alerting
dashboard into a single deployable pipeline that closes the loop from raw
Kraken2 reports to an actionable surveillance decision, eliminating the
analyst bottleneck described in Section~\ref{sec:intro}.

Within this pipeline, SBM community detection functions primarily as a
\emph{structural and interpretability} layer rather than an
independently validated risk signal: it consistently partitions the
co-occurrence graph into a biologically coherent
enteric/fecal-indicator cluster (\emph{Escherichia},
\emph{Enterobacter}, \emph{Citrobacter}, \emph{Vibrio};
Fig.~\ref{fig:sbm_network}) against a diffuse background. On this cohort the community-risk term saturates,
for a structural reason: under the $20\%$ pathogen-membership rule
\emph{both} $K=2$ communities qualify: the four-taxon enteric
cluster, and the 50-taxon background community, of which $24\%$
($12/50$) are database members because opportunistic wastewater
commensals (\emph{Pseudomonas}, \emph{Klebsiella}, \emph{Acinetobacter},
\emph{Streptococcus}, \emph{Enterococcus}) are themselves catalogued
pathogens. Since the background community spans most taxa in every
sample, every sample clears the $\phi=15\%$ gate and the term adds the
same $+0.25$ to every score, unable to discriminate \emph{between}
samples. Yet this uniform contribution is operationally consequential:
it is the deciding factor for 3 of the 5 real-data HIGH classifications
(Section~\ref{sec:real-data}), where abundance and novelty alone fall
just short of $\theta=0.6$. It becomes a per-sample \emph{discriminating}
signal only under a larger, sparser cohort or a membership rule
decoupled from community size.

\textbf{Limitations.} The real-data case study lacks clinical
ground-truth labels: the five HIGH-level runs and the elevated
novel-pathogen scores are plausible given the signal magnitudes but
cannot be confirmed as clinically relevant without paired
epidemiological records. The Spearman graph is built from $N=60$
samples, below the typical range for microbial co-occurrence
studies~\cite{Barber2012}; bootstrap analysis ($B=100$) keeps $K=2$ as
the modal partition (91\% of fits), but per-taxon co-assignment
probabilities within ($61.4\% \pm 39.7\%$) and between ($59.8\% \pm
26.6\%$) communities are not well separated, so the exact taxon-level
boundary is resampling-sensitive at this $N$. The saturated community
signal becomes a per-sample \emph{discriminating} signal only under the
conditions noted above (a larger, sparser cohort, or a size-decoupled
membership rule), which we leave to future work. Performance also
depends on Kraken2 database completeness, and the curated risk weights
require periodic expert review.

\textbf{Future work.} We plan to integrate the MaskedVQ-Seq
read-level novelty detector~\cite{Azumah2026} as an upstream signal,
extend the dashboard to multi-site geospatial visualization, and validate
on clinical metagenomic cohorts with confirmed pathogen outcomes.

In summary, PathogenIQ provides the public health and bioinformatics
community with a unified, open-source ML biosurveillance system operating
on standard Kraken2 outputs, producing interpretable risk scores, and
surfacing temporal outbreak signals in a real-time dashboard.

% ── References ────────────────────────────────────────────────────────────────
\begin{thebibliography}{00}

\bibitem{Bibby2021}
K.~J. Bibby et al., ``Making waves: Plausible lead time for wastewater based
epidemiology as an early warning system for COVID-19,''
\textit{Water Research}, vol.~202, p.~117438, 2021.

\bibitem{Wu2020}
F.~Wu et al., ``SARS-CoV-2 titers in wastewater are higher than expected
from clinically confirmed cases,'' \textit{mSystems}, vol.~5, no.~4, 2020.

\bibitem{Peccia2020}
J.~Peccia et al., ``Measurement of SARS-CoV-2 RNA in wastewater tracks
community infection dynamics,''
\textit{Nature Biotechnology}, vol.~38, pp.~1164--1167, 2020.

\bibitem{Medema2020}
G.~Medema et al., ``Presence of SARS-Coronavirus-2 RNA in sewage and
correlation with reported COVID-19 prevalence in the early stage of the
epidemic in The Netherlands,''
\textit{Environmental Science \& Technology Letters}, vol.~7, pp.~511--516, 2020.

\bibitem{Daughton2020}
C.~G. Daughton, ``Wastewater surveillance for population-wide COVID-19:
The present and future,''
\textit{Science of the Total Environment}, vol.~736, p.~139631, 2020.

\bibitem{CritsChristoph2021}
A.~Crits-Christoph et al., ``Genome sequencing of sewage detects regionally
prevalent SARS-CoV-2 variants,'' \textit{mBio}, vol.~12, e02703-20, 2021.

\bibitem{Hendriksen2019}
R.~S. Hendriksen et al., ``Global monitoring of antimicrobial resistance
based on metagenomics analyses of urban sewage,''
\textit{Nature Communications}, vol.~10, p.~1124, 2019.

\bibitem{Quince2017}
C.~Quince et al., ``Shotgun metagenomics, from sampling to analysis,''
\textit{Nature Biotechnology}, vol.~35, pp.~833--844, 2017.

\bibitem{Wood2019}
D.~E. Wood, J.~Lu, and B.~Langmead, ``Improved metagenomic analysis with
Kraken 2,'' \textit{Genome Biology}, vol.~20, p.~257, 2019.

\bibitem{Lu2017}
J.~Lu et al., ``Bracken: Estimating species abundance in metagenomics data,''
\textit{PeerJ Computer Science}, vol.~3, e104, 2017.

\bibitem{Breitwieser2019}
F.~P. Breitwieser, J.~Lu, and S.~L. Salzberg, ``A review of methods and
databases for metagenomic classification and assembly,''
\textit{Briefings in Bioinformatics}, vol.~20, pp.~1125--1136, 2019.

\bibitem{Liang2020}
Q.~Liang et al., ``DeepMicrobes: taxonomic classification for metagenomics
with deep learning,''
\textit{NAR Genomics and Bioinformatics}, vol.~2, lqaa009, 2020.

\bibitem{Azumah2026}
A.~Chinda, R.~Azumah, A.~Mikler, H.~Wang, and H.~Venkateswara,
``Reference-Free Variant Detection in Wastewater Genomic Surveillance
via Masked Vector-Quantized Autoencoders,''
\textit{Machine Learning for Healthcare}, 2026, under review.

\bibitem{Farrington1996}
C.~P. Farrington et al., ``A statistical algorithm for the early detection
of outbreaks of infectious disease,''
\textit{Journal of the Royal Statistical Society A}, vol.~159, pp.~547--563, 1996.

\bibitem{Unkel2012}
S.~Unkel et al., ``Statistical methods for the prospective detection of
infectious disease outbreaks: a review,''
\textit{Journal of the Royal Statistical Society A}, vol.~175, pp.~49--82, 2012.

\bibitem{Chandola2009}
V.~Chandola, A.~Banerjee, and V.~Kumar, ``Anomaly detection: A survey,''
\textit{ACM Computing Surveys}, vol.~41, no.~3, article~15, 2009.

\bibitem{Karthikeyan2022}
S.~Karthikeyan et al., ``Wastewater sequencing reveals early cryptic
SARS-CoV-2 variant transmission,'' \textit{Nature}, vol.~609, pp.~101--108, 2022.

\bibitem{Jahn2022}
K.~Jahn et al., ``Early detection and surveillance of SARS-CoV-2 genomic
variants in wastewater using COJAC,''
\textit{Nature Microbiology}, vol.~7, pp.~1785--1795, 2022.

\bibitem{Ahmed2020}
W.~Ahmed et al., ``First confirmed detection of SARS-CoV-2 in untreated
wastewater in Australia,''
\textit{Science of the Total Environment}, vol.~728, p.~138764, 2020.

\bibitem{Msomi2025}
N.~Msomi et al., ``Wastewater-integrated pathogen surveillance dashboards
enable real-time, transparent, and interpretable public health risk
assessment and dissemination,'' \textit{PLOS Global Public Health}, vol.~5,
no.~5, e0004443, 2025.

\bibitem{Gauthier2025}
J.~Gauthier et al., ``Leveraging artificial intelligence community analytics
and nanopore metagenomic surveillance to monitor early enteropathogen
outbreaks,'' \textit{Frontiers in Public Health}, vol.~13, art.~1675080, 2025.

\bibitem{Kohle2024}
S.~Kohle et al., ``Metagenomic analysis of sewage for surveillance of
bacterial pathogens: A release experiment to determine sensitivity,''
\textit{PLOS ONE}, vol.~19, no.~5, e0300733, 2024.

\bibitem{Feistel2025}
D.~J. Feistel et al., ``Detection and tracking of SARS-CoV-2 lineages through
national wastewater surveillance system pathogen genomics,''
\textit{Emerging Infectious Diseases}, vol.~31, suppl., pp.~57--60, 2025.

\bibitem{Pagsuyoin2025}
S.~Pagsuyoin et al., ``Coupling wastewater-based epidemiology with
data-driven machine learning for managing public health risks,''
\textit{Risk Analysis}, vol.~45, pp.~2974--2982, 2025.

\bibitem{Page1954}
E.~S. Page, ``Continuous inspection schemes,''
\textit{Biometrika}, vol.~41, pp.~100--115, 1954.

\bibitem{Holland1983}
P.~W. Holland, K.~B. Laskey, and S.~Leinhardt, ``Stochastic blockmodels:
First steps,'' \textit{Social Networks}, vol.~5, pp.~109--137, 1983.

\bibitem{Karrer2011}
B.~Karrer and M.~E.~J. Newman, ``Stochastic blockmodels and community
structure in networks,''
\textit{Physical Review E}, vol.~83, p.~016107, 2011.

\bibitem{Newman2006}
M.~E.~J. Newman, ``Modularity and community structure in networks,''
\textit{PNAS}, vol.~103, pp.~8577--8582, 2006.

\bibitem{Blondel2008}
V.~D. Blondel et al., ``Fast unfolding of communities in large networks,''
\textit{Journal of Statistical Mechanics}, p.~P10008, 2008.

\bibitem{Barber2012}
A.~Barber\'{a}n et al., ``Using network analysis to explore co-occurrence
patterns in soil microbial communities,''
\textit{ISME Journal}, vol.~6, pp.~343--351, 2012.

\bibitem{Kurtz2015}
Z.~D. Kurtz et al., ``Sparse and compositionally robust inference of
microbial ecological networks,''
\textit{PLOS Computational Biology}, vol.~11, e1004226, 2015.

\bibitem{Faust2012}
K.~Faust and J.~Raes, ``Microbial interactions: from networks to models,''
\textit{Nature Reviews Microbiology}, vol.~10, pp.~538--550, 2012.

\bibitem{Hadfield2018}
J.~Hadfield et al., ``Nextstrain: real-time tracking of pathogen evolution,''
\textit{Bioinformatics}, vol.~34, pp.~4121--4123, 2018.

\bibitem{Benjamini1995}
Y.~Benjamini and Y.~Hochberg, ``Controlling the false discovery rate,''
\textit{Journal of the Royal Statistical Society B}, vol.~57, pp.~289--300, 1995.

\bibitem{Peixoto2014}
T.~P. Peixoto, ``Efficient Monte Carlo and greedy heuristic for the
inference of stochastic block models,''
\textit{Physical Review E}, vol.~89, p.~012804, 2014.

\bibitem{WHO2024}
World Health Organization, ``WHO Bacterial Priority Pathogens List, 2024,''
WHO, Geneva, 2024.

\bibitem{Liu2008}
F.~T. Liu, K.~M. Ting, and Z.-H. Zhou, ``Isolation forest,''
in \textit{Proc. IEEE ICDM}, pp.~413--422, 2008.

\bibitem{Mann1945}
H.~B. Mann, ``Nonparametric tests against trend,''
\textit{Econometrica}, vol.~13, pp.~245--259, 1945.

\bibitem{Kendall1975}
M.~G. Kendall, \textit{Rank Correlation Methods}, 4th ed.
Charles Griffin, London, 1975.

\bibitem{Sen1968}
P.~K. Sen, ``Estimates of the regression coefficient based on Kendall's
tau,'' \textit{Journal of the American Statistical Association}, vol.~63,
pp.~1379--1389, 1968.

\end{thebibliography}

\end{document}
